\scanpage{33}\mspage{32}

\unsure{:} valable pour la ``topologie combinatoire sans points''. Donc il \unsure{faux|faut} revenir à la donnée de
\diagram{\begin{tikzcd}[ampersand replacement=\&,column sep=0.3em,row sep=1.1em]
C=\bigl(M, \& {\ll,}\arrow[d,out=-90,in=160] \& {R_M,}\arrow[d,out=-90,in=160] \& \text{Fam. de parties \deleted{\unsure{fermées}} de }M\bigr) \\
\& \substack{\text{raffinement}\\\text{(rel d'ordre}} \& \substack{\text{compatibilité}\\\text{rel, sym. réfl.}} \&
\end{tikzcd}}

% slanted
\margin{$R_M$ résulte de \unsure{``}$\ll,\mathcal F$, car \unsure{$\in$|si} \deleted{$F,G$} si $X,Y\in M$, $X\,R_M\,Y$ sss $\deleted{(X,Y)}$ $\exists F\in\mathcal F$ telle que $X,Y\in F$ D'autre part, on verra plus bas que si $C$ est finitaire i.e. les figures sont finies \inserted{(\unsure{cas|les} les plus importants !)}, alors $\ll,R_M$ déterminent $\mathcal F$.}

Je voudrais, pour commencer, au lieu de dégager un bon syst. d'axiomes, dégager une terminologie géométrique, en notant (s'il y a lieu) au passage les axiomes qui se seront imposés à l'attention en cours de route.

% slanted
\margin{\uline{NB} Une sous-contrée d'une sous-contrée est une \unsure{d|of} sous-contrée.}

\uline{Sous-contrée} : est définie par un sous-ensemble $M'$ \deleted{\uline{stable} \unsure{p|pr} $\ll$} \inserted{qui est} fermé dans $M$ pour $\ll$, i.e. si $X$ est dans $M'$ tout raffinement de $X$ y est aussi. \deleted{\unsure{Il|Je|Et}} \deleted{\unsure{Elle a} Mais on fera attention que} \inserted{On le munira des relations $\ll_{M'}$ et $R_{M'}$} induites par $\ll_M$ et $R_M$. Mais on fera attention qu'il n'est pas \deleted{évident} \inserted{dit} que les axiomes qu'on va imposer à une \unsure{contrée} soient \unsure{tous} stables par passage à une
